\begin{description}
\item[(H)] Areographic (Martian) latitude of the observer:
  $\varphi = 22.5^\circ$. \hfill(2 points)
\item[(I)] Pollux ($\beta$ Gem):
  $h_u = 61^\circ \pm 0.5^\circ$;
  $h_l = -50^\circ \pm 1^\circ$ (calculated). \hfill(1+2 points)\\
  Deneb ($\alpha$ Cyg):
  $h_u = 32^\circ \pm 0.5^\circ$;
  $h_l = 13^\circ \pm 0.5^\circ$. \hfill(1+1 points)
\item[(J)] Regulus ($\alpha$ Leo):
  $\delta = -22.5^\circ \pm 0.5^\circ$, from culmination
  $h_u = 45^\circ \pm 0.5^\circ$; \hfill(2 points)\\
  Toliman ($\alpha$ Cen):
  $\delta = -48^\circ \pm 0.5^\circ$, from culmination
  $h_u = 19.5^\circ \pm 0.5^\circ$. \hfill(2 points)
\item[(K)] Sketch diagrams on the meridian plane (zenith--nadir circle with
  the S--N horizon line): the celestial pole is drawn at altitude $\varphi$
  above the north point; the star's declination $\delta$ is laid off from
  the celestial equator to give the culmination altitude $h$.
  \hfill(2+2 points)
  % FIGURE NEEDED: two worked meridian-plane diagrams, 2011_PLAN_sol.pdf
  % page 3 (circles with Z, E, S, N, angles phi, delta, h_g marked)
\item[(L)] The Martian celestial North Pole is marked ``M'' on the answer
  map (in Cygnus, near the Cyg/Cep border). \hfill(2 points)
  % FIGURE NEEDED: marked answer sky map, 2011_PLAN_sol.pdf page 2
\item[(M)] $A = 330^\circ - 180^\circ = 150^\circ$. \hfill(2 points)
\item[(N)] (b) near the northern Tropic circle. \hfill(3 points)
\item[(O)] The date is marked on the time axis at the beginning of northern
  spring. \hfill(3 points)
  % FIGURE NEEDED (optional): season time axis with star marked at early
  % spring, 2011_PLAN_sol.pdf page 3 bottom
\end{description}
